现在假设 $v\in L^2(\Omega)$ 满足
\begin{equation}\label{eq:adjoint-weak-orthogonality-to-elliptic-range}
\int_\Omega(Lu)v\,\mathrm{d}\bm{x}=0
\end{equation}
对每个在 $\partial\Omega$ 上对 $j=1,\ldots,m$ 满足 $B_ju=0$ 的 $u\in H^{2m}(\Omega)$ 都成立. 若已知 $v\in H^{2m}(\Omega)$, 则可利用 Green 公式得到
\begin{equation}\label{eq:adjoint-weak-orthogonality-green-consequence}
\int_\Omega u(L^*v)\,\mathrm{d}\bm{x}=-\sum_{j=1}^m\int_{\partial\Omega}S_juC_jv\,\mathrm{d}S.
\end{equation}
由于 $B_ju$ 和 $S_ju$ 可以任意且彼此独立地选取, \eqref{eq:adjoint-weak-orthogonality-green-consequence} 意味着 $L^*v=0$, 且在边界上 $C_jv=0$. 即使不假设 $v\in H^{2m}(\Omega)$, 将 $u$ 限制于 $\mathcal{D}(\Omega)$ 后, 仍得到 $L^*v=0$ 在分布意义下成立. 剩下需要证明的是, 任意满足 \eqref{eq:adjoint-weak-orthogonality-to-elliptic-range} 的 $v\in L^2(\Omega)$ 实际上都属于 $H^{2m}(\Omega)$.

令
\begin{equation}\label{eq:adjoint-homogeneous-boundary-nullspace}
N^*=\{v\in H^{2m}(\Omega)\mid L^*v=0,\ C_1v=\cdots=C_mv=0\text{ on }\partial\Omega\},
\end{equation}
用 $(N^*)^\perp$ 表示 $N^*$ 在 $L^2(\Omega)$ 中的正交补, 并令 $M^*=(N^*)^\perp\cap H^{2m}(\Omega)$. 在 $M^*$ 上考虑二次型
\begin{equation}\label{eq:adjoint-auxiliary-quadratic-form}
[u,v]=\int_\Omega L^*uL^*v\,\mathrm{d}\bm{x}+\sum_{j=1}^m(C_ju,C_jv)_{2m-m_j-1/2},
\end{equation}
其中 $m_j$ 为 $C_j$ 的阶数, $(\cdot,\cdot)_s$ 表示 $H^s(\partial\Omega)$ 中的内积. 将证明这个二次型在空间 $M^*$ 上强制.

\begin{Lemma}
\label{lem:adjoint-auxiliary-form-complement-coercivity}
存在常数 $C$, 使得对所有 $u\in M^*$ 都有
\begin{equation}\label{eq:adjoint-auxiliary-form-complement-coercivity}
\|u\|_{2m}^2\leq C[u,u].
\end{equation}

\end{Lemma}

\begin{Proof}
由于边界算子 $C_j$ 对 $L^*$ 满足补条件, 可利用 Agmon--Douglis--Nirenberg 定理得到
\begin{equation}\label{eq:adjoint-auxiliary-form-garding-estimate}
\|u\|_{2m}^2\leq C([u,u]+\|u\|_0^2)
\end{equation}
对任意 $u\in H^{2m}(\Omega)$ 成立.

现在为导出矛盾, 假设 $u_n\in M^*$, $[u_n,u_n]\to0$, 而 $\|u_n\|_{2m}=1$. 则 $u_n$ 的某个子序列在 $H^{2m}(\Omega)$ 中弱收敛, 并在 $L^2(\Omega)$ 中强收敛; 设 $u$ 为极限. 由 \eqref{eq:adjoint-auxiliary-form-garding-estimate}, 该子序列实际上在 $H^{2m}(\Omega)$ 中强收敛到 $u$. 由于 $[u,u]=0$, 有 $u\in N^*$; 又由于 $u_n\in M^*$, 还有 $u\in M^*$. 因为 $N^*$ 和 $M^*$ 在 $L^2(\Omega)$ 中正交, 所以 $u=0$. 但这与 $\|u\|_{2m}=1$ 矛盾.
\end{Proof}

现在由 Lax--Milgram 引理, 方程 $[u,v]=(f,v)$ 对所有 $v\in M^*$ 成立时, 有唯一解 $u\in M^*$. 若 $f\in(N^*)^\perp$, 则实际上对所有 $v\in H^{2m}(\Omega)$ 都有 $[u,v]=(f,v)$. 可将其概括如下.
